\documentclass[11pt]{article}
\usepackage{amsmath,amssymb,amsthm,hyperref}
\newtheorem{theorem}{Theorem}
\newtheorem{lemma}{Lemma}
\newcommand{\E}{\mathbb E}
\newcommand{\Pol}{\operatorname{Pol}}
\title{A quadratic precision cutoff for approximately fair dice}
\author{Research proof; three independent full mathematical reviews completed}
\date{30 September 2026}
\begin{document}\maketitle
\noindent\textbf{Status.} Three independent full mathematical reviews completed
with no gap identified; this is not formal verification.
All constants below are absolute. Write $\log(1/0)=+\infty$.

\begin{theorem}\label{main}
Suppose $n\ge3$ independent equiprobable finite dice have no cross-die
ties, and every strict order has probability in
$[(1-\varepsilon)/n!,(1+\varepsilon)/n!]$, where
$0\le\varepsilon\le1/2$. Every die has at least
\[
 c n\min\{n,\log(1/\varepsilon)\}
\]
faces, for an absolute $c>0$.
\end{theorem}

We prove the following tensor version. Let $p_1,\ldots,p_m:[K]\to[0,1]$
be nondecreasing, with positive probability row weights $w$, and suppose
all complete failure patterns satisfy
\begin{equation}\label{approx}
 \E_w\prod_{j\in F}(1-p_j)\prod_{j\notin F}p_j
 \in\left[\frac{1-\varepsilon}{(m+1)\binom m{|F|}},
           \frac{1+\varepsilon}{(m+1)\binom m{|F|}}\right].
\end{equation}
Then both endpoint weights are at most
$C/[m\min\{m,H\}]$, where $H=\log(1/\varepsilon)$.
For dice, the event specifying which dice exceed the distinguished die
is a union of strict orders, so \eqref{approx} follows by summation.
Equal row weights then prove Theorem~\ref{main}.

Put $u_j=1-p_j$, $s=\sum_j u_j$, $\ell=\lfloor m/4\rfloor$,
$N=m-\ell$, $\alpha=\ell/m$, $x=\alpha s$, and $\theta=1/16$.
For a uniform $\ell$-subset $B$, let
$A_B=\prod_{j\in B}p_j$, $\bar A=\E_B A_B$, and
$v_B=N^{-1}\sum_{j\notin B}\ell u_j$.
The positive measure $\mu=(\ell+1)w\bar A$ has mass in
$[1-\varepsilon,1+\varepsilon]$; it is not normalized to probability.
Let $\nu=mw$ on the coordinate $x$, and let $\beta$ have density
\[
 b_\ell(z)=\frac{\ell+1}{\ell}(1-z/\ell)^\ell
                    \boldsymbol1_{[0,\ell]}(z).
\]

\begin{lemma}[Robust finite balance]\label{balance}
There is an absolute $\varepsilon_0>0$ such that, for
$\varepsilon\le\varepsilon_0$, the finite endpoint inputs hold with
absolute constants:
\begin{align}
 \max_j u_j&\le C(s+1)/m,\label{rough}\\
 \Pr_w(s\le S)&\le C(S+1)/m\quad(S\ge0),\label{mass}\\
 \E_w(s+1)^k e^{-a s}&\le C_a^{k+1}k!/m
                 \quad(k\ge0,a>0).\label{moments}
\end{align}
On $s\le c_0m$, all $p_j\ge1/2$, for an absolute $c_0>0$ and
sufficiently large $m$.
\end{lemma}
\begin{proof}
Summing \eqref{approx} over unspecified coordinates shows that every
partial pattern with $a$ successes and $b$ failures has probability
within relative $\varepsilon$ of
$\int_0^1t^a(1-t)^b\,dt$. This observation will also be used below.

Here are the quantitative changes to
\path{finite_endpoint_odds_balance.tex}. Normalize each $\ell$-failure
tilt by its actual mass, which is in $[1-\varepsilon,1+\varepsilon]$.
All its undivided generating coefficients, including the coefficients
with one specified derivative, then lie between the exact values
times $(1-\varepsilon)/(1+\varepsilon)$ and their reciprocals.
At a target row with $s\le m/256$, fix $i$ and choose
$J\subseteq[m]\setminus\{i\}$, $|J|=\lceil8s\rceil+4$, containing
all other columns with $p_j\le1/2$. Omit $i$ from the additional
failure-generating factors, as in that proof. The remaining target
odds sum is at most $2s\le|J|/4$.
For $\varepsilon_0$ sufficiently small, the same coefficientwise
bad-prefix bound with odds threshold one gives, for the good-row
remaining odds variable $Z$,
\[
 \E Z\ge\tfrac34(|J|+1),\qquad
 \E Z^2\le2(|J|+1)(|J|+3).
\]
Cauchy--Schwarz therefore gives probability at least $1/32$ on rows
strictly before the target. The specified odds expectation is at most
$2(|J|+1)/(m-|J|)$. As before, a zero target $p_i$ is impossible on
that prefix, and monotonicity yields $r_i\le C(s+1)/m$.
For larger $s$, (\ref{rough}) follows from $u_i\le1$.
Choosing a smaller absolute $c_0$ gives $p_j\ge1/2$ on $s\le c_0m$.

For (\ref{mass}), on $s\le S\le c_0m$ choose
$k=\lfloor c m/(S+1)\rfloor$ with a sufficiently small absolute $c$.
Every $k$-column success product is bounded below absolutely on that
suffix, while its expectation is at most $(1+\varepsilon)/(k+1)$.
This proves the mass bound; larger $S$ follow from total mass one.
Laplace integration and a supremum bound for $(s+1)^k e^{-as/2}$
prove (\ref{moments}). All constants are uniform in $\varepsilon$.
\end{proof}

\begin{lemma}[The source error has no dimension factor]\label{source}
For a polynomial $P(z)=\sum_{k=0}^D a_kz^k$, $D\le N$,
\[
 \left|(\ell+1)\E_w\E_B[A_B\Pol_N(P)((\ell u_j)_{j\notin B})]
                         -\beta(P)\right|
 \le\varepsilon\sum_k|a_k|k!.                  \tag{3}
\]
For the analogous marked measure with one fixed failure outside $B$,
normalization $(\ell+1)(\ell+2)$ and scalar density proportional to
$z(1-z/\ell)^\ell$, the bound is
$\varepsilon\sum_k|a_k|(k+1)!$.
If $P=Q_1Q_2$ and the degree-$d$ sums
$Q_b=\sum_{j\le d}c_{bj}L_j(\theta z)$ have coefficient sums $B_b$,
these errors are at most $\varepsilon B_1B_2 C_1^d$, for an absolute
$C_1>1$.
\end{lemma}
\begin{proof}
A degree-$k$ polarized monomial uses $k$ distinct failures disjoint
from the $\ell$ forced successes. Its exact normalized moment is
\[
 (\ell+1)\ell^k\int_0^1t^\ell(1-t)^k\,dt\le k!.
\]
The marked version has one more failure and its moment is at most
$(k+1)!$. The partial-pattern observation in Lemma~\ref{balance}
gives the relative error term by term, proving (3).
Finally, if $\|P\|_*=\sum|a_k|k!$, the explicit Laguerre coefficients
give
\[
 \|L_j(\theta z)L_k(\theta z)\|_*
 \le\sum_{a,b}\binom ja\binom kb2^{a+b}\le3^{j+k},
\]
since $(a+b)!/(a!b!)\le2^{a+b}$ and $\theta\le1$.
The marked norm costs only an additional factor $2d+1$, absorbed
by increasing $C_1$. Constants and degree zero are also covered by
increasing the prefactor when necessary.
\end{proof}

\begin{lemma}[Robust Gaussian comparison]\label{comparison}
For $\varepsilon\le\varepsilon_0$, $d\le c_1m$ with an absolute
$c_1>0$, and the Laguerre products of Lemma~\ref{source},
\[
 |\mu(Q_1Q_2)-\beta(Q_1Q_2)|
 \le C B_1B_2\{d/m+e^{-cm}+\varepsilon C_1^d\}.   \tag{4}
\]
The same assertion holds for the marked comparison after deletion of
one column. All constants are uniform in that column.
\end{lemma}
\begin{proof}
The analytic comparison in \path{../../power_per_die_lower_bound.tex}
uses only the finite bounds (\ref{rough})--(\ref{moments}), the
algebraic success-weighted subset identities, Maclaurin's inequality,
and the Gaussian Laguerre derivative and centered-variance estimates.
All these inputs remain valid. Split the subset comparison at
$s=c_0m$ rather than its former numerical cutoff. Its polarization
series is bounded by
$CB_1B_2\sum_{k\ge2}(C\sqrt{d/m})^k(k+1)^C$, so taking $c_1$
sufficiently small gives $CB_1B_2d/m$ even through linear degrees.
Lemma~\ref{source} adds the displayed source error.

For the marked measure, deletion leaves an approximate tensor of
$m-1$ columns by the partial-pattern observation. The original rough
bound rearranges to $u_i\le C(s_{-i}+1)/m$. Thus its normalization
$O(m^2)$ gives weighted factorial moments $C^{k+2}(k+1)!$, rather
than $C^{k+1}k!$. This adds only a polynomial factor to the same
geometric polarization series. The marked source error is precisely
the second assertion of Lemma~\ref{source}. No probability
normalization of $\mu$ is needed in this analytic argument.
\end{proof}

We retain the endpoint tests
\[
 T_d(y)=\sum_{j<d}(1-j/d)L_j(y),\quad
 U_d(z)=d^{-2}T_d(\theta z)^2,\quad
 R_d(z)=(2/(\theta d)-z)T_d(\theta z)^2.
\]
Their reviewed scalar properties are
\[
 \beta(R_d)\ge c>0,\quad \beta(U_d)\le C/d,
 \quad U_d(z)\ge c\quad(0\le z\le c_2/d).
 \tag{5}
\]
The coefficient product norm is $O(1)$ for $U_d$ and $O(d)$ for
$R_d$, using $yT_d=d^{-1}\sum_{j<d}L_j-L_d$.

\begin{lemma}[High-precision inputs]\label{high}
Suppose $H\ge\sqrt m$. Then, for sufficiently large $m$,
\begin{align}
 \max_j u_j&\le C(s+\delta)/m,&\delta&=m^{-1/3},\label{improved}\\
 \nu([0,t])&\le C(t+h),&h&=\max\{H^{-1},m^{-3/4}\}.\label{cdf}
\end{align}
For $\varepsilon=0$, take $H^{-1}=0$.
\end{lemma}
\begin{proof}
First (4), applied to $U_d$ with $d\le c_3\sqrt m$, gives
$\mu(U_d)\le C/d+Cd/m+\varepsilon C_1^d$. Choose $c_3>0$
small enough that $\varepsilon C_1^d\le e^{-H/2}$. Near-zero
positivity in (5), $\bar A\ge1-s$, and degree selection therefore
give the initial CDF bound $\nu([0,t])\le C(t+m^{-1/2})$.

For the relative balance, repeat the marked Laplace argument in
\path{quantitative_relative_endpoint_balance.tex}. Truncate the
geometric Laguerre expansion of $e^{-tz/2}$ at
$D=\lfloor c_3\sqrt m\rfloor$, reducing $c_3$ further if needed.
Its coefficients $c_j=(1-r)r^j$, $r=t/(t+2\theta)$, have sum at
most one and weighted degree sum $O(t)$. Applying (4) entrywise to
$L_jL_k$ therefore gives, for either marked or unmarked measure,
\[
 |\mu_a(e^{-t x})-\beta_a(e^{-tz})|
 \le Ct/m+C e^{-c\sqrt m/t}+C e^{-H/2}.
 \tag{6}
\]
This is uniform for $1/2\le t\le m^{1/3}$. The source errors were
bounded by their maximum $\varepsilon C_1^D$, while the analytic
errors use the weighted degree sum. Both measures may have masses
$1\pm\varepsilon$; positivity and the displayed estimates suffice.
Their Laplace bounds and the unmarked CDF bound give fixed constants
$a,b,c,C$ with
$\mu_0(a/t<x<b/t)\ge c/t$ and
$\mu_1(x\le b/t)\le C/t^2$, exactly as in that note.
The unchanged monotone density relation
$d\mu_1=(\ell+2)u_i\,d\mu_0$ proves (\ref{improved}).

For (\ref{cdf}), use the elementary all-orders $U_d$ comparison in
\path{../../size_bounds/elementary_endpoint_cdf_bootstrap.tex}.
Its analytic proof uses only the rough balance and initial CDF;
by Lemma~\ref{source} the approximate version is
\[
 |\mu(U_d)-\beta(U_d)|
 \le C\{m^{-1}+m^{-1/2}d/m+\log d/(dm)\}
       +O(m^{-3})+\varepsilon C_1^d.
 \tag{7}
\]
This holds for $m^{1/4}\le d\le c_4\min(H,m^{3/4})$, with
$c_4$ a sufficiently small absolute constant. The source term is
at most $e^{-H/2}$, and the right side is $O(1/d)$.
The usual near-zero positivity and degree selection give
$\nu([0,t])\le C(t+1/D_{max})$,
$D_{max}=c_4\min(H,m^{3/4})$. This is (\ref{cdf}).
\end{proof}

\begin{proof}[Proof of the tensor bound and Theorem~\ref{main}]
The previously reviewed approximate linear lower bound, proved in
\path{approximate_per_die_lower_bound.tex}, covers bounded $H$ with
endpoint weights $O(1/m)$, uniformly for $\varepsilon\le1/2$.
Thus assume $H\ge H_0$, with $H_0$ a sufficiently large absolute
constant, so $\varepsilon\le\varepsilon_0$. Small bounded $m$ are
covered by enlarging constants.

If $H\le\sqrt m$, choose $d=\lfloor c_5H\rfloor$. Equations
(4)--(5) bound the $R_d$ error by
\[
 C d^2/m+C d e^{-cm}+C\varepsilon d C_1^d.
\]
First choose $c_5$ small, then $H_0$ large. The first term is at most
$Cc_5^2$, and the source is at most $Cd e^{-H/2}$, so the error is
less than half the positive scalar expectation. Therefore
$x_K<2/(\theta d)$. Apply $U_r$, $r=\lfloor d/16\rfloor$;
it is bounded below at $x_K$, and its comparison error is at most
$Cr/m+Ce^{-cm}+C e^{-H/2}=O(1/d)$. Hence $\mu(K)\le C/d$.
Since $s_K=O(1/d)$, $\bar A(K)\ge1/2$ after increasing $H_0$,
and $w_K\le C/(md)\le C/(mH)$.

Now let $H\ge\sqrt m$, and use Lemma~\ref{high}. Choose
$d=\lfloor c_6\min(H,m)\rfloor$, with a sufficiently small fixed
$c_6>0$. We record explicitly the uniform estimates obtained from
\path{../../size_bounds/quadratic_per_die_lower_bound.tex} for this
possibly sublinear degree. Its all-orders Gaussian argument, using
(\ref{improved})--(\ref{cdf}), gives
\begin{align*}
 |\mu(R_d)-\beta(R_d)|
 &\le C\frac d m(1+\delta\log(2d)+\delta hd)
       +o(1)+C\varepsilon d C_1^d,\\
 |\mu(U_d)-\beta(U_d)|
 &\le C\left\{\frac{\log(2d)}{dm}+\frac hm
                   +\frac\delta m+\frac{\delta hd}m\right\}
       +o(m^{-2})+C\varepsilon C_1^d.             \tag{8}
\end{align*}
Here the little-oh tails are uniform for
$c\sqrt m\le d\le c_6m$. To verify this without relying on a
fixed proportional degree, the small-$x$ errors are respectively
$C(d^2/m)(1/d+\delta)(1/d+h)$ and
$C(d/m)(1/d+\delta)(1/d+h)$, absorbed in (8).
The middle-region envelopes are
$C(d/m)e^{-cx}(1+\delta/x)$ and
$C(dm)^{-1}e^{-cx}(x^{-1}+\delta x^{-2})$.
Their CDF integrals are precisely the terms in (8).
Bad subsets have exponent at least $cm/(d\delta)\ge c m^{1/3}$;
the outer region has exponent at least $cd\ge c\sqrt m$.
Choose $c_6$ small enough for the same exponential margins as in
the exact proof. Lemma~\ref{source} supplies the only additional
terms, the displayed source errors.

In this regime
\[
 \delta hd\le C m^{-1/3}\max\{1,m^{1/4}\}=O(m^{-1/12}),
 \qquad \varepsilon C_1^d\le e^{-H/2}.
\]
Thus the first line of (8) is at most $Cc_6+o(1)$, made smaller
than half the scalar positivity constant. Again $x_K<2/(\theta d)$.
For $r=\lfloor d/16\rfloor$, the second line of (8) is $o(1/d)$:
after multiplication by $d$, every displayed term tends to zero.
Since $U_r(x_K)\ge c$ and $\beta(U_r)\le C/r$, we obtain
$\mu(K)\le C/d$. The same endpoint conversion gives
\[
 w_K\le C/(md)\le C/[m\min(H,m)].
\]
Reversing rows and complementing all columns preserves the approximate
pattern hypothesis and gives the same result for $w_1$.
The tensor-to-dice reduction at the start completes the proof.
\end{proof}
\end{document}
